% Overleaf main document: main.tex (slides); choose cover.tex for the A4 cover.
\documentclass[aspectratio=169,11pt]{beamer}
\usepackage[T1]{fontenc}
\usepackage{lmodern}
\usepackage{graphicx}
% Edit here: shared by the Beamer slides and the A4 publication cover.
% Use LaTeX escapes for special characters: \&, \%, \_, \#.
\newcommand{\DPTitle}{Title of the detector performance note}
\newcommand{\DPNumber}{2026/XXX} % Placeholder; use the assigned number.
\newcommand{\DPDate}{\today} % Or set a fixed date.
\newcommand{\DPAuthors}{CMS Collaboration} % A4 cover credit.
\newcommand{\DPContact}{cms-dpg-conveners-mtd@cern.ch}
\newcommand{\DPAbstract}{%
Replace this text with the final abstract. Explain the purpose of the study,
the detector or reconstruction method, the data or simulation used, and the
main performance results. Define abbreviations on first use and make the
summary understandable without consulting the plots.

This single abstract is used on slide 2 and on the A4 cover page.
Keep it short enough to fit on one slide; replace this guidance before publication.
}


% Plain, black-on-white Beamer theme.
\usetheme{default}
\setbeamercolor{normal text}{fg=black,bg=white}
\setbeamercolor{structure}{fg=black}
\setbeamercolor{frametitle}{fg=black,bg=white}
\setbeamerfont{frametitle}{size=\large,series=\bfseries}
\setbeamersize{text margin left=8mm,text margin right=8mm}
\setbeamertemplate{navigation symbols}{}
% A dedicated header reserves space for the logo on EVERY slide.
% Do not use [plain] frames: those suppress Beamer headers and footers.
\setbeamertemplate{headline}{%
  \hbox to \paperwidth{\hfill
    \raisebox{0pt}[12mm][1mm]{\includegraphics[height=10mm]{CMSlogo_black_label_1024.png}}%
    \hspace{6mm}}%
}
\setbeamertemplate{footline}{%
  \hbox to \paperwidth{\hspace{8mm}\scriptsize CMS DP-\DPNumber
    \hfill\insertframenumber\hspace{8mm}}\vspace{3mm}%
}
\setbeamertemplate{frametitle}{%
  \vspace{-1mm}\insertframetitle\par\vspace{2mm}%
}

% Each call keeps the required label attached to its plot panel.
% For plots already carrying this exact label, use \DPPlotEmbedded instead.
\newcommand{\DPPlotLabel}{{\small\textbf{CMS} Preliminary Phase 2\par}\vspace{1mm}}
\newcommand{\DPPlot}[2][43mm]{%
  \begin{minipage}{\linewidth}
    \DPPlotLabel
    \centering\includegraphics[width=\linewidth,height=#1,keepaspectratio]{#2}\par
  \end{minipage}%
}
\newcommand{\DPPlotEmbedded}[2][48mm]{%
  \centering\includegraphics[width=\linewidth,height=#1,keepaspectratio]{#2}\par
}
\newcommand{\DPPlotPlaceholder}{%
  \begin{minipage}{\linewidth}
    \DPPlotLabel
    \setlength{\fboxsep}{0pt}%
    \fbox{\begin{minipage}[c][41mm][c]{\dimexpr\linewidth-2\fboxrule\relax}
      \centering\color{black!55}Plot placeholder\par
      \small Replace with your measured or simulated result.
    \end{minipage}}%
  \end{minipage}%
}

\hypersetup{pdftitle={\DPTitle},pdfauthor={CMS}}
\begin{document}

% 1. Title and requested author/contact line.
\begin{frame}
  \centering
  {\LARGE\DPTitle\par}
  \vspace{9mm}
  {\large CMS\par}
  \vspace{3mm}
  \href{mailto:\DPContact}{\DPContact}
\end{frame}

% 2. Same abstract as the publication cover page.
\begin{frame}{Abstract}
  \small
  \setlength{\parskip}{5pt}
  \DPAbstract
\end{frame}

% 3. Duplicate this frame for each result and replace ALL placeholders.
% A caption must stand on its own: explain the sample, axes and units,
% selections, curves/markers, uncertainties, and what the result means.
\begin{frame}[t]{Descriptive title of the performance result}
  \begin{columns}[T,onlytextwidth]
    \begin{column}{0.56\textwidth}
      \DPPlotPlaceholder
      % Upload a PDF/PNG plot into figures/, then replace the line above with:
      % \DPPlot{figures/your_plot.pdf}
      % If the plot already has the exact required label, use:
      % \DPPlotEmbedded{figures/your_plot.pdf}
    \end{column}
    \begin{column}{0.41\textwidth}
      \fontsize{10}{12}\selectfont
      \textbf{Caption placeholder}\par\medskip
      This plot shows [quantity, meaning and units] versus [variable and
      units], using [sample] with [conditions and selection].
      \par\medskip
      The [markers/curves] represent [samples or methods]. The
      [error bars/bands] show [statistical and/or systematic uncertainties].
      \par\medskip
      The [trend] indicates [conclusion and significance for detector
      performance]. Define all abbreviations.
    \end{column}
  \end{columns}
\end{frame}

\end{document}
